\documentclass[10pt,aspectratio=169]{beamer}
\usepackage{bma01}

\title[Concavidad e inflexión]{Concavidad de gráficas}
\subtitle{Puntos de inflexión. Criterio de la segunda derivada\\[2pt]\small Semana 11 -- Sesión 1}

\begin{document}

\begin{frame}
  \titlepage
\end{frame}

\begin{frame}{Objetivos de la sesión}
  \begin{enumerate}
    \item Definir \resaltar{concavidad} y relacionarla con el signo de $f''$.
    \item Hallar \resaltar{puntos de inflexión}.
    \item Aplicar el \resaltar{criterio de la segunda derivada} para extremos locales.
    \item Completar el esbozo de gráficas con la información de $f''$.
  \end{enumerate}
  \begin{block}{Plan (2 horas)}
    \begin{tabular}{ll}
      0:00 -- 0:30 & Concavidad: definición geométrica y prueba con $f''$\\
      0:30 -- 0:55 & Puntos de inflexión. Aplicaciones\\
      0:55 -- 1:15 & Criterio de la segunda derivada\\
      1:15 -- 1:45 & Esbozo completo de gráficas\\
      1:45 -- 2:00 & Ejercicios y resumen
    \end{tabular}
  \end{block}
\end{frame}

\section{Concavidad}

\begin{frame}{Concavidad}
  \begin{columns}
    \begin{column}{0.5\textwidth}
      \begin{definicion}
        Sea $f$ derivable en un intervalo $I$.
        \begin{itemize}
          \item $f$ es \resaltar{cóncava hacia arriba} (convexa) en $I$ si su gráfica está \textbf{sobre} todas sus tangentes en $I$
            (equivalentemente, $f'$ es creciente).
          \item $f$ es \resaltar{cóncava hacia abajo} en $I$ si su gráfica está \textbf{bajo} sus tangentes ($f'$ decreciente).
        \end{itemize}
      \end{definicion}
      \begin{teorema}[criterio de concavidad]
        Si $f''>0$ en $I$: cóncava hacia arriba. \ Si $f''<0$ en $I$: cóncava hacia abajo.
      \end{teorema}
    \end{column}
    \begin{column}{0.5\textwidth}
      \centering
      \begin{tikzpicture}
        \begin{axis}[bma, xmin=-2.2, xmax=2.2, ymin=-0.5, ymax=4, width=6.4cm, height=3.6cm, xtick=\empty, ytick=\empty, samples=100, title={\scriptsize hacia arriba ($\cup$): tangentes por debajo}]
          \addplot[curva, domain=-2:2] {x^2};
          \addplot[curva2, thin, domain=-1:2] {2*(x-1)+1};
          \addplot[curva2, thin, domain=-2:0.5] {-3*(x+1.5)+2.25};
        \end{axis}
      \end{tikzpicture}
      \begin{tikzpicture}
        \begin{axis}[bma, xmin=-2.2, xmax=2.2, ymin=-4, ymax=0.5, width=6.4cm, height=3.6cm, xtick=\empty, ytick=\empty, samples=100, title={\scriptsize hacia abajo ($\cap$): tangentes por encima}]
          \addplot[uniVerde, very thick, domain=-2:2] {-x^2};
          \addplot[curva2, thin, domain=-1:2] {-2*(x-1)-1};
          \addplot[curva2, thin, domain=-2:0.5] {3*(x+1.5)-2.25};
        \end{axis}
      \end{tikzpicture}
    \end{column}
  \end{columns}
\end{frame}

\begin{frame}{¿Por qué funciona el criterio?}
  \begin{itemize}
    \item Si $f''>0$, entonces $f'$ es creciente (TVM aplicado a $f'$).
    \item Para $x>a$: por el TVM, $f(x)-f(a)=f'(c)(x-a)$ con $a<c<x$, y $f'(c)>f'(a)$. Luego
      \[ f(x)>f(a)+f'(a)(x-a), \]
      es decir, la gráfica está sobre la tangente en $a$. (Análogo para $x<a$.)
  \end{itemize}
  \begin{idea}
    \begin{itemize}
      \item $f'$ dice si la gráfica \textbf{sube o baja}.
      \item $f''$ dice si la pendiente \textbf{aumenta o disminuye}: la ``curvatura'' de la gráfica.
    \end{itemize}
  \end{idea}
  \begin{ejemplo}[física]
    Si $s(t)$ es la posición, $s''=a$ es la aceleración: $s$ cóncava hacia arriba $\iff$ la velocidad aumenta.
  \end{ejemplo}
\end{frame}

\section{Puntos de inflexión}

\begin{frame}{Puntos de inflexión}
  \begin{columns}
    \begin{column}{0.55\textwidth}
      \begin{definicion}
        $P=(c,f(c))$ es un \resaltar{punto de inflexión} si $f$ es continua en $c$ y la concavidad \textbf{cambia} en $c$.
      \end{definicion}
      \begin{block}{Candidatos}
        Puntos donde $f''(c)=0$ o $f''(c)$ no existe. \textbf{Hay que verificar el cambio de signo de $f''$}.
      \end{block}
      \begin{ejemplo}[1]
        \begin{itemize}
          \item $f(x)=x^3$: $f''=6x$ cambia de signo en $0$: inflexión en $(0,0)$.
          \item $g(x)=x^4$: $g''=12x^2=0$ en $0$, pero no cambia de signo: \textbf{no} hay inflexión.
          \item $h(x)=\sqrt[3]x$: $h''$ no existe en $0$ y cambia de signo: inflexión (con tangente vertical).
        \end{itemize}
      \end{ejemplo}
    \end{column}
    \begin{column}{0.45\textwidth}
      \centering
      \begin{tikzpicture}
        \begin{axis}[bma, xmin=-1.6, xmax=1.6, ymin=-1.6, ymax=2.2, width=6.3cm, height=5.5cm, xtick=\empty, ytick=\empty, samples=100]
          \addplot[curva, domain=-1.2:1.2] {x^3};
          \addplot[curva2, domain=-1.2:1.2] {x^4};
          \addplot[only marks, mark=*, black] coordinates {(0,0)};
          \node[uniAzul] at (axis cs:1.35,1.4) {\scriptsize $x^3$};
          \node[uniGranate] at (axis cs:-1.4,1.9) {\scriptsize $x^4$};
        \end{axis}
      \end{tikzpicture}
    \end{column}
  \end{columns}
\end{frame}

\begin{frame}{Inflexión en las ciencias}
  \begin{columns}
    \begin{column}{0.5\textwidth}
      \begin{ejemplo}[2: campana de Gauss]
        $f(x)=e^{-x^2/2}$: \ $f''(x)=(x^2-1)e^{-x^2/2}$.
        Inflexiones en $x=\pm1$: en la distribución normal, a \textbf{una desviación estándar} de la media.
      \end{ejemplo}
      \centering
      \begin{tikzpicture}
        \begin{axis}[bma, xmin=-3.5, xmax=3.5, ymin=0, ymax=1.2, width=6cm, height=3.5cm, xtick={-1,1}, ytick=\empty, samples=100, domain=-3.4:3.4]
          \addplot[curva] {exp(-x^2/2)};
          \addplot[only marks, mark=*, uniGranate] coordinates {(-1,0.6065) (1,0.6065)};
        \end{axis}
      \end{tikzpicture}
    \end{column}
    \begin{column}{0.5\textwidth}
      \begin{ejemplo}[3: curva logística]
        $P(t)=\dfrac{K}{1+Ae^{-rt}}$ tiene inflexión cuando $P=\frac K2$: el crecimiento es \textbf{máximo} a mitad de la capacidad.
      \end{ejemplo}
      \begin{ejemplo}[4: titulación]
        En una curva de titulación pH vs.\ volumen, el \textbf{punto de equivalencia} es el punto de inflexión (máxima pendiente).
      \end{ejemplo}
      \centering
      \begin{tikzpicture}
        \begin{axis}[bma, xlabel={$V$}, ylabel={pH}, xmin=0, xmax=50, ymin=0, ymax=14, width=6cm, height=3.5cm, xtick={25}, ytick={7}, samples=200, domain=0:50]
          \addplot[curva2] {7+5.5*rad(atan((x-25)/1.2))/1.5708};
        \end{axis}
      \end{tikzpicture}
    \end{column}
  \end{columns}
\end{frame}

\section{Criterio de la segunda derivada}

\begin{frame}{Criterio de la segunda derivada}
  \begin{teorema}
    Si $f'(c)=0$ y $f''$ es continua cerca de $c$:
    \begin{itemize}
      \item $f''(c)>0$ $\Rightarrow$ \resaltar{mínimo local} en $c$ \ (el ``vaso'' $\cup$ retiene agua);
      \item $f''(c)<0$ $\Rightarrow$ \resaltar{máximo local} en $c$;
      \item $f''(c)=0$ $\Rightarrow$ el criterio \textbf{no decide} (use la primera derivada).
    \end{itemize}
  \end{teorema}
  \begin{ejemplo}[5]
    $f(x)=x^4-4x^3$: \ $f'(x)=4x^2(x-3)$, \ $f''(x)=12x^2-24x=12x(x-2)$.
    \begin{itemize}
      \item $x=3$: $f''(3)=36>0$: mínimo local $f(3)=-27$.
      \item $x=0$: $f''(0)=0$: no decide. Con $f'$: negativa a ambos lados $\Rightarrow$ no hay extremo.
      \item Inflexiones: $f''$ cambia de signo en $x=0$ y $x=2$: puntos $(0,0)$ y $(2,-16)$.
    \end{itemize}
  \end{ejemplo}
\end{frame}

\begin{frame}{Gráfica del ejemplo 5}
  \begin{columns}
    \begin{column}{0.5\textwidth}
      \centering
      \begin{tikzpicture}
        \begin{axis}[bma, xmin=-1.5, xmax=4.6, ymin=-32, ymax=15, width=7cm, height=6.4cm, xtick={2,3,4}, ytick={-16,-27}, samples=100, domain=-1.3:4.4]
          \addplot[curva] {x^4-4*x^3};
          \addplot[only marks, mark=*, uniGranate] coordinates {(3,-27)};
          \addplot[only marks, mark=*, uniVerde] coordinates {(0,0) (2,-16)};
        \end{axis}
      \end{tikzpicture}
    \end{column}
    \begin{column}{0.5\textwidth}
      \begin{center}
        \small
        \begin{tabular}{c|cccc}
          & $(-\infty,0)$ & $(0,2)$ & $(2,3)$ & $(3,\infty)$\\ \hline
          $f'$ & $-$ & $-$ & $-$ & $+$\\
          $f''$ & $+$ & $-$ & $+$ & $+$\\
          forma & $\searrow\cup$ & $\searrow\cap$ & $\searrow\cup$ & $\nearrow\cup$
        \end{tabular}
      \end{center}
      \begin{itemize}
        \item Granate: mínimo local (y absoluto).
        \item Verde: puntos de inflexión.
        \item En $(0,0)$ la tangente es horizontal, pero es inflexión, no extremo.
      \end{itemize}
    \end{column}
  \end{columns}
\end{frame}

\section{Esbozo completo}

\begin{frame}{Esbozo completo: $f(x)=\dfrac{x}{x^2+1}$}
  \begin{columns}
    \begin{column}{0.55\textwidth}
      \begin{itemize}
        \item Dominio $\R$; impar; corte $(0,0)$; asíntota horizontal $y=0$.
        \item $f'(x)=\dfrac{1-x^2}{(x^2+1)^2}$: crece en $(-1,1)$; máx.\ $f(1)=\frac12$, mín.\ $f(-1)=-\frac12$.
        \item $f''(x)=\dfrac{2x(x^2-3)}{(x^2+1)^3}$: cambia de signo en $x=0,\pm\sqrt3$.
          \begin{center}
            \small
            \begin{tabular}{c|cccc}
              & $(-\infty,-\sqrt3)$ & $(-\sqrt3,0)$ & $(0,\sqrt3)$ & $(\sqrt3,\infty)$\\ \hline
              $f''$ & $-$ & $+$ & $-$ & $+$
            \end{tabular}
          \end{center}
        \item Inflexiones: $(0,0)$, $\left(\pm\sqrt3,\pm\frac{\sqrt3}{4}\right)$.
      \end{itemize}
    \end{column}
    \begin{column}{0.45\textwidth}
      \centering
      \begin{tikzpicture}
        \begin{axis}[bma, xmin=-6, xmax=6, ymin=-0.7, ymax=0.7, width=6.6cm, height=5cm, xtick={-1.732,-1,1,1.732}, xticklabels={$-\sqrt3$,,,$\sqrt3$}, ytick={-0.5,0.5}, samples=200, domain=-6:6]
          \addplot[curva] {x/(x^2+1)};
          \addplot[only marks, mark=*, uniGranate] coordinates {(1,0.5) (-1,-0.5)};
          \addplot[only marks, mark=*, uniVerde] coordinates {(0,0) (1.732,0.433) (-1.732,-0.433)};
        \end{axis}
      \end{tikzpicture}
    \end{column}
  \end{columns}
\end{frame}

\section{Ejercicios}

\begin{frame}{Ejercicios propuestos}
  \small
  \begin{ejercicio}
    \begin{enumerate}
      \item Halle intervalos de concavidad y puntos de inflexión: (a) $f(x)=x^3-6x^2+9x$, \ (b) $g(x)=xe^{-x}$, \ (c) $h(x)=\ln(1+x^2)$,
        \ (d) $k(x)=x^{1/3}(x+4)$, \ (e) $p(x)=\sen x+\cos x$ en $[0,2\pi]$.
      \item Use el criterio de la segunda derivada: (a) $f(x)=x+\frac4x$, \ (b) $g(x)=x^2\ln x$, \ (c) $h(x)=\cos^2x$ en $[0,\pi]$.
      \item Halle $a,b$ para que $f(x)=ax^3+bx^2$ tenga un punto de inflexión en $(1,2)$.
      \item Pruebe que todo polinomio cúbico tiene exactamente un punto de inflexión, y que la gráfica es simétrica respecto de él.
      \item Esboce con toda la información: (a) $f(x)=\dfrac{x^2}{x^2+3}$, \ (b) $g(x)=\dfrac{x^3}{x^2-1}$, \ (c) $h(x)=x^2e^{-x}$.
      \item Pruebe que la logística $P(t)=\dfrac{K}{1+Ae^{-rt}}$ tiene su inflexión en $P=\frac K2$.
      \item Dibuje una función continua en $\R$ con $f'>0$, $f''<0$ para $x<0$ y $f'>0$, $f''>0$ para $x>0$. ¿Qué tiene en $x=0$?
    \end{enumerate}
  \end{ejercicio}
\end{frame}

\begin{frame}{Resumen}
  \begin{itemize}
    \item $f''>0$: cóncava hacia arriba ($\cup$), gráfica sobre las tangentes; $f''<0$: hacia abajo ($\cap$).
    \item Inflexión: cambio de concavidad. Candidatos: $f''=0$ o $f''$ no existe; verificar el cambio de signo.
    \item Segunda derivada: $f'(c)=0$ y $f''(c)>0$ $\Rightarrow$ mínimo; $f''(c)<0$ $\Rightarrow$ máximo; $=0$: no decide.
    \item Esbozo completo: dominio, simetría, cortes, asíntotas, $f'$ (monotonía, extremos), $f''$ (concavidad, inflexiones).
  \end{itemize}
  \vspace{3mm}
  \textbf{Próxima sesión:} la regla de L'Hospital.
\end{frame}

\end{document}
